\documentclass[11pt,a4paper]{article}
\usepackage[margin=0.82in]{geometry}
\usepackage[T1]{fontenc}
\usepackage{lmodern,microtype}
\usepackage{amsmath,amssymb,amsthm,mathtools}
\usepackage{booktabs,tabularx,array,enumitem}
\usepackage[hidelinks]{hyperref}

\newtheorem{definition}{Definition}[section]
\newtheorem{lemma}[definition]{Lemma}
\newtheorem{theorem}[definition]{Theorem}
\newtheorem{proposition}[definition]{Proposition}
\newtheorem{corollary}[definition]{Corollary}
\newtheorem{guard}[definition]{Guard}
\newcommand{\PASS}{\texttt{PASS}}
\newcommand{\REF}{\texttt{REFUTED}}
\newcommand{\SEP}{\texttt{SEPARATED}}
\newcommand{\NT}{\texttt{NT}}

\title{\textbf{Pointed Involutions, KEP \(30\mid14\) Normal Forms, and the Roof--Rail Reflection Frontier}\\
\large A proof-first KEP/RH/Stone refinement of the \(44\)-split and the \(89=44+45\) closure}
\author{Formal working note}
\date{27 August 2026}

\begin{document}
\maketitle

\begin{abstract}
We refine the roof--rail recurrence comparison by introducing a finite
\(C_2\)-set invariant and by replaying source-bound Kepler (KEP), RH/A1BE,
and Stone carriers.  An exact A1BE pointed \(89\)-element involution is
constructed from the \(44\) complete even reversal two-cycles together
with the distinguished fixed point \((3,26)\).  It has one fixed point,
\(44\) free two-cycles, and \(45\) total orbits; consequently
\[
89=44+45
\]
is realized structurally, not merely arithmetically.

KEP-R13 supplies the native decomposition
\[
44=9+21+7+7,\qquad 30_{\rm even}=9+21,\quad14_{\rm odd}=7+7.
\]
KEP-R27 supplies the exact Cube27 atomic decomposition
\[
27=1+12+6+8=(1+12)+(6+8)=13+14.
\]
These data give canonical transfer sizes \(21,13,12\) for the three
recurrence-positive roof branches \(9\mid35\), \(17\mid27\), and
\(18\mid26\).  The negative branch \(8\mid36\) would require transfer
size \(22\), which is not a union of the declared R13 face classes nor
of the declared R27 atomic blocks.  This is a scoped no-go only:
the independent Stone width-three reversal has an exact \(22\)-state
partial-action domain, with escape labels \(8,17\).

The RH rectangle of \(36\) points has six fixed points and fifteen
two-cycles; its involution split is \(36=15+21\), supplying an independent
structural realization of the positive mover \(21\).  The remaining
roof/KEP/RH-to-recurrence commuting square is left explicitly
\(\SEP/\NT\).
\end{abstract}

\section{Finite involutions and the canonical orbit split}

\begin{definition}
Let \(X\) be a finite set equipped with an involution \(J\).
Write
\[
f=|\operatorname{Fix}(J)|,\qquad
p=\#\{\text{free two-cycles of }J\},
\qquad
q=|X/J|.
\]
\end{definition}

\begin{theorem}[Involution split]
For every finite involution,
\[
|X|=2p+f,\qquad q=p+f,
\]
and therefore
\[
\boxed{|X|=p+q.}
\]
Moreover \(q-p=f\).
\end{theorem}

\begin{proof}
The \(J\)-orbits are precisely \(f\) singleton fixed orbits and \(p\)
two-element free orbits.  Hence \(|X|=f+2p\) and \(q=f+p\).
\end{proof}

\begin{corollary}[Pointed case]
If \(J\) has a unique fixed point, then
\[
|X|=2p+1,\qquad |X/J|=p+1,
\]
so
\[
\boxed{|X|=p+(p+1).}
\]
\end{corollary}

\section{An exact RH/A1BE pointed \(89\)-carrier}

Let \(E\) be the \(212\)-element even support in the frozen A1BE source
and let
\[
J(i,j)=(29-j,29-i).
\]
The source replay has exactly seven fixed support pairs and \(44\)
complete two-cycles.  Let \(T_{88}\) be the union of the \(88\) points
in those \(44\) complete two-cycles and let
\[
p_\star=(3,26)
\]
be the F1-distinguished fixed pair.  Define
\[
P_{89}=T_{88}\sqcup\{p_\star\}.
\]

\begin{theorem}[Source-bound pointed \(89\)-involution]
The set \(P_{89}\) is \(J\)-stable, has cardinality \(89\), and has
exactly one fixed point.  Consequently
\[
\boxed{
\#\{\text{free }J\text{-orbits}\}=44,\qquad
|P_{89}/J|=45,
}
\]
and hence
\[
\boxed{89=44+45}
\]
as the free-orbit count plus the total-orbit count of one explicit
source-bound involution.
\end{theorem}

\begin{proof}
The \(44\) complete two-cycles contribute \(88\) pairwise distinct
points and are \(J\)-stable.  The distinguished point \(p_\star\) is
fixed and is not in a two-cycle.  Thus \(P_{89}\) has one fixed orbit
and \(44\) free orbits.  Apply the pointed involution corollary.
\end{proof}

\begin{guard}
This theorem does not identify the A1BE \(P_{89}\) carrier with the
prime \(89\), with the Fibonacci carrier \(F_{11}\), or with the
order-\(89\) recurrence carrier.  It realizes only the same integer
decomposition \(89=44+45\) by a genuine involutive structure.
\end{guard}

\section{The RH rectangle \(36\) and a native \(15\mid21\) split}

The F1 rectangle component has \(36=6\cdot6\) points and is closed under
the same reversal.  Its exact orbit census is
\[
f=6,\qquad p=15.
\]
Hence
\[
q=p+f=21
\]
and
\[
\boxed{36=15+21.}
\]
Equivalently, for the permutation representation,
\[
\dim V_+=21,\qquad \dim V_-=15.
\]
The number \(21\) therefore occurs as an RH-native quotient/even
dimension, independently of the KEP \(F4opp\) count used below.

\section{KEP-R13: a native \(9\mid35\to30\mid14\) normal form}

The recovered R13 endpoint carrier has four face classes:
\[
|F6|=9,\quad |F4opp|=21,\quad |F5opp|=7,\quad |F5adj|=7.
\]
Grouping by active-edge parity yields
\[
\boxed{
30_{\rm even}=9+21,\qquad
14_{\rm odd}=7+7.
}
\]

Therefore the split
\[
9\mid35
=
9\mid(21+7+7)
\]
has the native KEP transfer
\[
\boxed{
9\mid35
\xrightarrow{\;F4opp_{21}\;}
(9+21)\mid(7+7)
=
30\mid14.
}
\]

This is a genuine decomposition of one R13 object carrier, not merely a
comparison of dimensions.

\section{KEP-R27: Cube27 transfer blocks}

Inside the \(27\)-point carrier \(\{-1,0,1\}^3\), the declared atomic
classes have cardinalities
\[
\boxed{
1\quad\text{(centre)},\qquad
12\quad\text{(cuboctahedral roots)},\qquad
6\quad\text{(axes)},\qquad
8\quad\text{(cube corners)}.
}
\]
The parity split is
\[
\boxed{
C_{13}=1+12,\qquad O_{14}=6+8,\qquad27=13+14.
}
\]

Consequently, after explicitly declaring a tagged right-hand Cube27
carrier,
\[
\boxed{
17\mid27
=
17\mid(13+14)
\longrightarrow
(17+13)\mid14
=
30\mid14.
}
\]
Likewise, on the punctured Cube27 carrier
\[
26=12+6+8=12+14,
\]
so
\[
\boxed{
18\mid26
=
18\mid(12+14)
\longrightarrow
(18+12)\mid14
=
30\mid14.
}
\]

\begin{guard}
Unlike the R13 \(9\mid35\) transfer, the preceding two displays are
typed tagged-carrier normal forms.  They use native KEP subcarriers on
the right, but do not assert that the left \(17\)- or \(18\)-object
roof/recurrent carrier is literally a KEP subset.
\end{guard}

\section{The scoped \(22\)-mover obstruction}

For a split \(a\mid(44-a)\), transfer to \(30\mid14\) by moving a
right-hand block to the left requires a block of size
\[
m(a)=30-a.
\]
For the four roof-selected branches:
\[
\begin{array}{c|c|c}
\text{split}&m(a)&\text{recurrence reflection}\\\hline
8\mid36&22&\text{no}\\
9\mid35&21&\text{yes}\\
17\mid27&13&\text{yes}\\
18\mid26&12&\text{yes}
\end{array}
\]

Let \(\mathcal M_{\rm R13}\) be the subset-sum set generated by
\(9,21,7,7\), and \(\mathcal M_{\rm R27}\) the subset-sum set generated
by \(1,12,6,8\).  Then
\[
21\in\mathcal M_{\rm R13},\qquad
13,12\in\mathcal M_{\rm R27},
\]
but
\[
\boxed{
22\notin\mathcal M_{\rm R13}\cup\mathcal M_{\rm R27}.
}
\]

Thus the three recurrence-positive roof branches have a transfer by one
of the declared canonical KEP block unions, whereas \(8\mid36\) does
not.

\begin{guard}[Finite alignment, not a general classifier]
This agreement is proved only on the four roof-selected branches.
The KEP subset-sum criterion does not reproduce the complete
\(21\)-pair recurrence classification at \(k=44\), and it is not
promoted to a causal explanation of recurrence reflection.
\end{guard}

\section{Stone supplies the missing cardinal \(22\), but on another carrier}

Encode \(r\in\{0,\ldots,23\}\) as a width-three ternary word and reverse
its digits.  The closed partial-action domain is
\[
\boxed{
D_{22}=\{0,\ldots,23\}\setminus\{8,17\}.
}
\]
Indeed
\[
8=022_3\mapsto220_3=24,\qquad
17=122_3\mapsto221_3=25.
\]
Thus \(|D_{22}|=22\).  On \(D_{22}\), reversal has eight fixed points
and seven free two-cycles.  On the minimal total closure
\(\{0,\ldots,25\}\), it has eight fixed points and nine two-cycles,
so the involution split gives
\[
\boxed{26=9+17.}
\]

Therefore \(22\) is not globally absent.  The precise current statement
is:
\[
\boxed{
22\text{ is absent from the declared canonical KEP mover class,}
}
\]
while an independent Stone \(22\)-carrier exists and has the rail
escape labels \(8,17\).  Constructing a structure-preserving map from
this Stone carrier into the \(8\mid36\) normal-form problem is a new,
sharply typed target.

\section{F1 Fourier \(9/21\) and the correct categorical level}

On the \(6\times6\) Fourier grid, independent antipody invariance leaves
nine modes.  Under the swap involution \(J(p,q)=(q,p)\), the
\(J\)-even sector has
\[
6+15=21
\]
modes and the \(J\)-odd sector has \(15\) modes.  Hence the
\(9/21/15\) profile is exact.

Later KEP continuation constructs deterministic \(9\to F6\) and
\(21\to F4opp\) representative maps and then upgrades them to
orbit-valued functors.  The native orbit saturations have sizes
\(60\) and \(42\), so a direct equivariant bijection of the original
unsaturated sets is impossible.  This identifies the appropriate
future category: orbit/groupoid-valued transport rather than raw
finite-set identity.

\section{A prediction for \(34\mid55\)}

For an \(89\)-point involution,
\[
p=\frac{89-f}{2},\qquad
q=\frac{89+f}{2}.
\]
Therefore
\[
\boxed{
p=34,\quad q=55
\quad\Longleftrightarrow\quad
f=21.
}
\]

Thus a structural realization of
\[
89=34+55
\]
by the same involution-split mechanism would require a native
\(89\)-point \(C_2\)-carrier with exactly \(21\) fixed points.

No such source-bound carrier is established in the present replay.
The statement is therefore an \(\NT\) target, not a theorem of
existence.

\section{Remaining recurrence commuting square}

The recurrence reflection acts on roots by
\[
\rho_{\rm rec}(\alpha)=-\alpha^{-1}.
\]
The next proof target is an explicit typed map \(\Psi\) from one of the
source involution/groupoid carriers above into the recurrence
factor-orbit carrier such that
\[
\boxed{
\Psi\circ J_{\rm source}
=
\rho_{\rm rec}\circ\Psi.
}
\]
At present:
\[
\begin{array}{p{0.64\linewidth}p{0.22\linewidth}}
\toprule
\textbf{Statement}&\textbf{Status}\\\midrule
A1BE pointed $P_{89}$: \(44\) free cycles and \(45\) total orbits
&\PASS\\
Structural involution identity \(89=44+45\)
&\PASS\\
RH rectangle \(36=15+21\)
&\PASS\\
Native KEP \(9\mid35\to30\mid14\)
&\PASS\\
Tagged Cube27 \(17\mid27\to30\mid14\) via \(13\)
&\PASS\ typed\\
Tagged punctured-Cube27 \(18\mid26\to30\mid14\) via \(12\)
&\PASS\ typed\\
Canonical KEP block-union mover \(22\) for \(8\mid36\)
&\REF\ scoped\\
Stone \(22\)-domain with escape labels \(8,17\)
&\PASS\ separate\\
Native \(89\)-point involution with \(21\) fixed points
&\NT\\
Full source-to-recurrence commuting square
&\SEP/\NT\\
\bottomrule
\end{array}
\]

\section*{Source locks}

The replay used the following source SHA-256 values:
\begin{itemize}[leftmargin=2em]
\item \texttt{A1BE\_direct\_tail\_pairs.jsonl}:\\
\texttt{7fa2ac47faa32aa5170da5e18bb4fff236239b9d52f55021292b75a3e907e881};
\item \texttt{A1BE\_F1\_reversal\_analytic\_morphism.json}:\\
\texttt{b3ff6920ef9b1303892f5ced198fc1f33154194682891a9856bb545ea436a8c6};
\item \texttt{II\_IO\_KEP\_R13\_stable44\_ledger.json}:\\
\texttt{62b9b76f3aec2a4e6102f83be144cca94ba25e1856885e2d9278c77b422843df};
\item \texttt{kep\_d3\_fcc\_box\_transport\_formalization\_R27.tex}:\\
\texttt{1a95d8f8c4dcdfdf16a6c605d1effa40b0ab4b7c17a201290212d1de0a1312c4};
\item \texttt{II\_HW\_partialgroupoid\_zRquotient\_ReesAutomaton\_H3cocycle\_R1.tex}:\\
\texttt{449c3e82733abe02fa2a439867865b17cbc03cc8e91bab66bc9f0deb12138fe6}.
\end{itemize}

\end{document}
